\documentclass[11pt]{amsart}
\usepackage[a4paper,margin=1in]{geometry}
\usepackage[T1]{fontenc}
\usepackage{lmodern,microtype,amsmath,amssymb,booktabs,array}
\usepackage[hidelinks]{hyperref}
\newtheorem{theorem}{Theorem}[section]
\newtheorem{lemma}[theorem]{Lemma}
\newtheorem{proposition}[theorem]{Proposition}
\theoremstyle{remark}\newtheorem{remark}[theorem]{Remark}
\newcommand{\R}{\mathbb R}
\newcommand{\one}{\mathbf1}
\newcommand{\T}{\mathsf T}
\newcommand{\tr}{\operatorname{tr}}
\newcommand{\diag}{\operatorname{diag}}
\title[Constructive Koml\'os discrepancy: final certificate]{Constructive Koml\'os discrepancy below 37.54:\newline final proof and algorithm specification}
\author{}
\date{28 September 2026}
\begin{document}
\begin{abstract}
We refine the project's retained-row spectral signing walk to obtain the
exact constructive discrepancy bound $37.54$. The improvement comes from
counting spectral cancellation and positive curvature together, then
normalizing the response to account jointly for the large-row deposit and protected row constraints.
The implementation uses ordinary field operations and comparisons,
including rational approximations rather than square-root primitives.
Its proved total cost remains $O(mn)+\widetilde O(n^4)$. A randomized
variant reduces all row processing to $O(mn)+\widetilde O(n^3)$ expected
work, but construction of the changing signing subspaces still costs
quartic work. No cubic total-time or polynomial bit-complexity claim is made.
\end{abstract}
\maketitle

\section{Result, scope, and the walk being modified}

\begin{theorem}\label{thm:main41}
Let $A\in\R^{m\times n}$ have columns of Euclidean norm at most one.
There is a deterministic algorithm returning one sign per original column,
$\varepsilon\in\{-1,1\}^n$, such that
\[
 \|A\varepsilon\|_\infty<\frac{1877}{50}=37.54.
\]
It uses $O(n\log^4(2+n))$ updates and
$O(mn)+\widetilde O(n^4)$ ordinary field operations and comparisons.
No fast matrix-multiplication exponent is used. Rational input gives
rational states; a polynomial bound for their bit lengths is not proved.
\end{theorem}

This is a technical addendum to the below-$67$ manuscript~\cite{baseline},
whose complete source and certificates are included in this repository. We retain its cleanup, tracking, rational Jacobi interface,
exact projection, flattening, screening, and termination arguments.
The response and dimension estimates, their finite-step compatibility,
and every changed scalar parameter are proved below. Every changed scalar parameter is
listed below, and the supporting core source and checks are included in
the same repository. Thus the result is
an arithmetic-operation theorem, not a benchmarked numerical solver.

The retained-row framework comes from the project draft
\cite{project}. Its screening and motion-clock ingredients draw on
Li~\cite{Li}. The contribution here is the joint inertia calculation, its normalization,
and the common dimension budget for protected pairs, large rows, and
spectral response. The randomized variance clock in
Section~\ref{sec:randomscreen} is a separate implementation refinement.
We make no priority claim beyond this derivation.

\begin{center}
\begin{tabular}{lcc}
\toprule
Guarantee & Previous joint walk & This addendum\\
\midrule
Discrepancy & $40.88$ & $37.54$\\
Ordinary arithmetic, total & $O(mn)+\widetilde O(n^4)$
 & $O(mn)+\widetilde O(n^4)$\\
Uniform cubic total bound & unproved & unproved\\
\bottomrule
\end{tabular}
\end{center}

\subsection{State and derivative notation}
Start at $x=0$ in the cube, with active coordinates $V$ and $s=|V|$.
Row $i$ has retained coefficients $b_{ij}$, offset $c_i$, squared vector
$p_{ij}=b_{ij}^2$, mass $q_i=\sum_jp_{ij}$, tracked sum
$d_i=c_i+b_i^Tx_V$, and energy $E_i=\sum_jp_{ij}f(x_j)$.
Large rows have $q_i>R^2$ and impose $b_i^Th=0$.
Medium rows have $0<q_i\le R^2$. Delete their entries satisfying
$p_{ij}>\Delta q_i$, transferring $b_{ij}x_j$ to the offset.
This preserves $d_i$ and makes the row diffuse. Use the same nonincreasing
rational scales $q_i\le r_i^2<\theta^2q_i$, $r_i\le R$ as in
\cite{baseline}.

For row-sign pairs $r=(i,\sigma)$ write
\[
 u_r=\frac{\sigma d_i-H_0}{r_i}+\frac{E_i}{r_i^2}+b,
 \qquad a_r=A_0e^{\beta u_r},\qquad z_r=p_i/r_i^2.
\]
The monitored matrix and potential are
\[
 W=\sum_r a_rz_rz_r^T,\quad
 B=B_0\sum_{i\text{ large}}\diag(p_i),\quad
 M=W+B+\eta\one\one^T,\quad
 \Phi=(\lambda_1(M)-1)_+^3-\eta\|x\|^2.
\]
The norm of $x$ includes frozen coordinates. During differentiation,
retained data and scales are held fixed. Put
\[
 g_r=\sigma b_i/r_i+z_r\circ f'(x),\quad
 \mathcal U=\sum_r a_r\beta^2(z_r^Tv)^2g_rg_r^T,\quad
 \mathcal D=\sum_r a_r\beta(z_r^Tv)^2
                      \diag(z_r\circ(-f''(x))),
\]
where $v$ is the positive unit Perron vector at the anchor.
The profile below gives $\mathcal U_{jj}\le\rho\mathcal D_{jj}$.
These formulas also hold for the frozen-cold comparison curve, with
derivatives taken only over warm pairs.

At phase start $S$, set $J=1+\lceil\log_2 n\rceil$,
$\eta=1/(2^{40}JS)$ and $\delta=\eta^5/2^{100}$; reset when
$s\le S/2$. Discard input rows of $\ell_1$ norm at most 37 in one scan.
The surviving count satisfies $m_0<n^2/1369$, and
$q=\lceil\log_2(8(n+2m_0))\rceil=O(\log(2+n))$.
In the following lemmas $\ell_0=999999/1000000$ is the fixed lower bound
for $L$ in the high branch. In the low branch $L\le1$, no normalized response is needed.
The elementary bound on large and protected rows leaves at least
$[1-R^{-2}-c_pL_*-2/337-1/4094]s>s/400$ feasible dimensions;
the exact parameter verifier also checks this separate branch.

\section{Joint inertia of spectral cancellation and coupled response}

This section replaces the separate charges for high spectral directions
and positive directions of $\gamma\mathcal U-\mathcal D$.
It uses the actual low-spectrum response. All computations below use
ordinary matrix multiplication and elimination.

\begin{lemma}[Reciprocal Gram inertia]
Let $D\succeq0$ be diagonal, let $G\in\R^{r\times s}$, and suppose
$(G^TG)_{jj}\le\rho D_{jj}$. Let $F\in\R^{r\times r}$ satisfy $F\succeq I$,
and write its eigenvalues decreasingly as $f_1,f_2,\ldots$, padding by ones
if necessary. If $G^TFG-D$ has $k$ positive eigenvalues, then
\[
 \rho s\ge\operatorname{tr}(D^{\dagger/2}G^TG D^{\dagger/2})
 >\sum_{j=1}^k f_j^{-1}.
\]
The strict comparison is asserted for $k>0$; for $k=0$ it is replaced
by a weak comparison.
\end{lemma}
\begin{proof}
A zero diagonal entry of $D$ forces the corresponding column of $G$ to
vanish. Thus both matrices vanish on that coordinate and it can be removed.
Congruence by $D^{-1/2}$ preserves inertia. Put $Y=GD^{-1/2}$ and choose an
orthonormal matrix $X$ spanning the positive eigenspace of $Y^TFY-I$.
Then $B^2=X^TY^TFYX\succ I$. The columns of
$H=F^{1/2}YX B^{-1}$ are orthonormal, and
\[
 \operatorname{tr}(Y^TY)\ge\|YX\|_F^2
 =\operatorname{tr}(B^2 H^TF^{-1}H)
 >\operatorname{tr}(H^TF^{-1}H)
 \ge\sum_{j=1}^k f_j^{-1}.
\]
The last inequality is the variational principle for the sum of the
smallest eigenvalues. The diagonal hypothesis bounds the first trace.
\end{proof}

\begin{lemma}[A rational joint count]\label{lem:newjoint}
Suppose $N=VDV^T\succeq0$, where $V$ is exactly orthogonal, $D$ is diagonal,
$\lambda_* =\max_jD_{jj}>0$, and
$CC^T\preceq N+\epsilon I$. Let $F_0$ be the columns with
$D_{jj}\ge(18/25)\lambda_*$ and let $H=|F_0|$. Put
\[
 R=V_{F_0^c}\operatorname{diag}((\lambda_*-D_{jj})^{-1})V_{F_0^c}^T,
 \qquad Q=G^TG-D_0+2G^TC^TRCG,
\]
where $D_0\succeq0$ is diagonal and $(G^TG)_{jj}\le\rho(D_0)_{jj}$.
If $Q$ has $k$ positive eigenvalues, then
\[
 H+k < {25\over23}\rho s+{32\over23}{\operatorname{tr}N\over\lambda_*}
            +{50\over23}{\epsilon s\over\lambda_*}.
\]
The weak inequality version applies when $k=0$.
\end{lemma}
\begin{proof}
Apply the preceding lemma to $F=I+2C^TRC$. The nonzero eigenvalues of
$R^{1/2}CC^TR^{1/2}$ are bounded, in decreasing order, by the numbers
$(D_{jj}+\epsilon)/(\lambda_*-D_{jj})$ on the low space.
This follows from the stated PSD domination and the variational principle.
Consequently each relevant reciprocal eigenvalue of $F$ is at least
\[
 {1-\lambda\over1+\lambda+2\epsilon/\lambda_*}
 \ge {1-\lambda\over1+\lambda}-{2\epsilon\over\lambda_*}
 \ge {23\over25}-{32\over25}\lambda-{2\epsilon\over\lambda_*},
 \qquad \lambda=D_{jj}/\lambda_*.
\]
The last inequality is the supporting tangent at $\lambda=1/4$;
its difference from the left-hand rational function is exactly
$32(\lambda-1/4)^2/[25(1+\lambda)]\ge0$.
Pad the low list by zeros if $k$ is larger than its length. Summing gives
\[
 {23\over25}k
 <\rho s+{32\over25}{\operatorname{tr}N_{\rm low}\over\lambda_*}
       +{2\epsilon s\over\lambda_*}.
\]
Since $\lambda\ge18/25>23/32$ on the high list,
$(23/25)H\le(32/25)\operatorname{tr}N_{\rm high}/\lambda_*$.
Add these inequalities. The $k=0$ case follows just from the high-list
bound and nonnegativity of all the other terms.
\end{proof}

\section{Normalizing away the large-row deposit}

Write $M=B+W+\eta\one\one^T$, where $B$ is the nonnegative diagonal
large-row deposit. At an anchor with top eigenvalue $L>\ell_0$ put
\[
 D_L=LI-B,\qquad S_L=D_L^{-1/2},\qquad
 K_L=S_L(W+\eta\one\one^T)S_L.
\]
Throughout this section $B_0\le7/8$, so $D_L\succ0$.
The matrix $K_L$ is PSD and entrywise positive, has Perron eigenvalue one,
and has positive unit Perron vector
$u=D_L^{1/2}v/t_0$, where $t_0=\|D_L^{1/2}v\|$.
Its gap below one is at least $\eta/L>\eta/2$.
Indeed $LI-M=D_L^{1/2}(I-K_L)D_L^{1/2}$, and the nonzero eigenvalues
of a PSD matrix under a congruence are bounded using the minimum singular
value of the congruence. This gives the asserted gap from that of $M$.

If $y\perp v$, then
\begin{equation}\label{eq:genresponse}
 y^T(LI-M)^\dagger y
 =(S_Ly)^T(I-K_L)^\dagger(S_Ly).
\end{equation}
To check this without a pseudoinverse congruence rule, note that
$S_Ly\perp u$. The vector
$S_L(I-K_L)^\dagger S_Ly$ solves $(LI-M)z=y$; it differs from the reduced
solution only by a multiple of $v$, whose inner product with $y$ is zero.
No assertion that pseudoinverses themselves commute with congruence is used.

Include the ridge as an additional Gram column in $C$, and append a zero
row to $G$ for that column. Include every frozen cold Gram term in the
same fashion. The derivative forcing is $CG$; its varying rows still
give $G^TG=\mathcal U$ and diagonal negative term $\mathcal D$.
The normalized Gram matrix is $C_LC_L^T=K_L$, with $C_L=S_LC$.
Thus the reciprocal Gram count applies to the normalized low response.

For each warm row-sign pair, approximate its barrier with absolute error at
most $1/1000$, and protect it when the computed value is at least $-1/200$.
Impose the exact equation $g_r^Th=0$ for every protected pair. Thus a
protected pair has true $u_r\ge-1/100$, and an unprotected warm pair has
$u_r<-1/250$. The next lemma counts these equations jointly with the
large-row and spectral constraints.

\begin{lemma}[Joint charge for protected pairs, large rows, and response]
Use a reciprocal-Gram supporting line $\varphi(\lambda)\ge A_t-B_t\lambda$
with $A_t>0$, and write $a_t=1/A_t$, $b_t=B_t/A_t$.
Let
\[
 c_p=\frac{\theta^4}{A_0(1-\beta/100)}.
\]
Suppose the witness cutoff is at least $A_t/B_t$ and
\begin{equation}\label{eq:depositcondition}
 B_0R^2\left(c_p+\frac{b_t\Delta}{\ell_0-B_0}\right)\le1.
\end{equation}
Apart from the already budgeted precision and ridge errors, the total
number of large-row equations, protected-gradient equations, high
normalized spectral equations, and positive curvature directions is
at most
\[
 \left(\frac1{R^2}+a_t\rho+b_t\Delta+c_p(L_*-B_0)\right)s.
\]
\end{lemma}
\begin{proof}
A protected pair has $u\ge-1/100$ and $\sum_jz_j=q/r^2>\theta^{-2}$.
It therefore contributes at least
$A_0\theta^{-4}(1-\beta/100)=1/c_p$ to $\one^TW\one$.
If $p$ is the protected-pair count and $\ell$ is the large-row count,
\[
 p\le c_p( Ls-\operatorname{tr}B-\eta s^2),\qquad
 \operatorname{tr}B\ge B_0R^2\ell.
\]
The normalized trace and reciprocal-Gram count give
\[
 H+k\le a_t\rho s+b_t
       \frac{\Delta(Ls-\operatorname{tr}B)+\eta s}{L-B_0}
\]
up to the diagonalization error. Add $p$ and $\ell$, discard the negative
ridge term, and use the trace lower bound. The coefficient of $\ell$ is
\[
 1-B_0R^2\left(c_p+\frac{b_t\Delta}{L-B_0}\right)\ge0.
\]
It is therefore valid to substitute $\ell\le s/R^2$. The two terms with
denominator $L-B_0$ cancel, and the protected charge becomes $c_p(L-B_0)$.
Finally use $L\le L_*$. The positive ridge remainder is bounded by
$b_t\eta s/(\ell_0-B_0)$, and fits the existing numerical allowance.
\end{proof}

For a tangent at $t\in[0,\sqrt2-1)$,
\[
 A_t=\frac{1+2t-t^2}{(1+t)^2},\qquad
 B_t=\frac2{(1+t)^2},\qquad
 a_t=\frac{(1+t)^2}{1+2t-t^2},\quad
 b_t=\frac2{1+2t-t^2}.
\]
The witness cutoff only needs to exceed
$A_t/B_t=(1+2t-t^2)/2$; the cutoff $18/25$ therefore permits
$t<1-\sqrt{14}/5$. The identity behind the tangent is
\[
 \frac{1-\lambda}{1+\lambda}-(A_t-B_t\lambda)
 =\frac{2(\lambda-t)^2}{(1+\lambda)(1+t)^2}.
\]
This changes the counting certificate only. The same response proxy,
constraints, norm bounds, finite-step lemma, and arithmetic model apply.

\paragraph{Explicit perturbation allowance.}
The existing numerical contracts give
$\|N_K-K_L\|\le256\delta$, $|k_*-1|\le256\delta$, and
$C'C'^T\preceq N_K+\delta I$. Since $\operatorname{tr}K_L\le s$,
\[
 \frac{\operatorname{tr}N_K}{k_*}
 \le\operatorname{tr}K_L+1024\delta s.
\]
For the relevant tangents $a_t,b_t\le2$, the resulting error in the
joint spectral count is at most $2^{12}\delta s$; this includes the
reciprocal-line perturbation $2a_t\delta s/k_*$.
The positive ridge remainder is at most
$b_t\eta s/(\ell_0-B_0)\le18\eta s$ when $B_0\le7/8$.
Thus the fully explicit additive allowance is
\[
 (18\eta+2^{12}\delta)s<\frac{s}{10000},
 \qquad \eta\le2^{-40},\quad\delta=\eta^5/2^{100}.
\]
All other terms of the dimension certificate, including the two exact
tangencies and the row-energy high-space exclusion, are unchanged.

\paragraph{Flat tests and cold rows.}
Use test vectors $e_j$ and $(1-c)g_{i\sigma}$, whose Euclidean norms are
at most one. This explicitly replaces the old profile's normalization
$(8/25)g_{i\sigma}$, which is not valid for every improved profile.
For $h=h_0/16$ the dyadic test based on $9q/k$ gives
$16\|h\|_\infty\le K$ and $|\beta g_{i\sigma}^Th|\le K$, provided
$\beta/(1-c)\le16$. The stronger checked bound
$\beta/(1-c)\le25/8$ also preserves the original numerical forcing
and frozen-cold estimates. In particular every cold row satisfies
$|\ell'(0)|\le25/128<1/4$; combined with the coordinate flatness and
second-derivative bound, this gives $|\ell(t)|<|t|/4$ on the prescribed
step interval. Replace every occurrence of the old amplitude in the
cold threshold by the current $A_0$: it suffices to choose the least
integer $\ell_c$ such that $2^{\ell_c}\ge A_0/\rho_c$.

The finite-step and proxy certificates only require
\[
 \frac{\beta c}{256d(1-c)^2}<\frac12,\qquad
 \frac{2\beta c^2}{4096d(1-c)^3}<\frac14,\qquad
 \frac{\beta cd}{(1-c)^2}\le\frac94,
 \qquad \frac\beta{1-c}\le\frac{25}{8}.
\]
These are rational checks at rational parameters.

\paragraph{Rational implementation of the normalized proxy.}
First compute the approximate top pair $(\lambda_*,w)$ of $M$ as in the
original walk, with $\|M-N_M\|<\delta$ and $\|w-v\|\le3\delta/\eta$
after orientation. Approximate the positive diagonal
$(\lambda_*I-B)^{-1/2}$ by a rational diagonal matrix $S$, with absolute
error at most $\delta/2^{20}$. Its entries lie in a fixed bounded interval;
bisection of positive square roots to this accuracy costs
$O(s\log(1/\delta))$ field operations and comparisons. Square roots are
not operation primitives. Form
\[
 \bar K=S(\bar W+\eta\one\one^T)S
\]
and apply rational Jacobi diagonalization at accuracy $\delta/4$ to get an exactly orthogonal
matrix $V_K$ and diagonal $D_K$, and clip negative reported diagonal entries to zero. The final reconstructed
matrix $N_K$ is within $\delta$ of $\bar K$.
Let $k_* =\max_j(D_K)_{jj}$. Put in the high set $F$ every reported column
with $(D_K)_{jj}\ge(18/25)k_*$. Set
\[
 R_K=(V_K)_{F^c}\operatorname{diag}((k_*-(D_K)_{jj})^{-1})(V_K)_{F^c}^T.
\]
The low denominators are bounded below by $(7/25)k_*$.
Compute, using only the warm pairs,
\[
 \bar U=\sum_r\bar a_r\beta^2(z_r^Tw)^2g_rg_r^T,\quad
 \bar D=\sum_r\bar a_r\beta(z_r^Tw)^2
                         \operatorname{diag}(z_r\circ(-f'')),
 \quad \bar B_1=\sum_r\bar a_r\beta(z_r^Tw)z_rg_r^T.
\]
Add the exact equations $(V_K)_F^T S\bar B_1h=0$ and use the actual proxy
\begin{equation}\label{eq:genproxy}
 \bar Q=\bar U-\bar D+2\bar B_1^T S R_K S\bar B_1.
\end{equation}
The diagonal inequality $\bar U_{jj}\le\rho\bar D_{jj}$ is exact.
For analysis, the normalized computed Gram factors satisfy
$C'C'^T=\bar K\preceq N_K+\delta I$, so Lemma~\ref{lem:newjoint}
applies directly to the computed proxy and computed equations.
This avoids using a discontinuous cutoff perturbation theorem.

\paragraph{Accuracy contracts.}
The scalar inverse-square-root derivative and $B_0\le7/8$ give
$\|S-S_L\|\le16\delta$, and $\|S\|,\|S_L\|<3$.
Thus $\|K_L-N_K\|\le256\delta$, $|k_*-1|\le256\delta$,
and the computed top column $w_K$ satisfies
$\|w_K-u\|\le2048\delta/\eta$.
The true normalized gap is at least $\eta/2$ and the reconstructed gap
is at least $\eta/4$.
The high normalized forcing obeys, for exactly feasible $h$,
\[
 \|P_{(3/4,1]}(K_L)S_LM'[h]v\|,
 \ |L'[h]|
 \le 2^{30}\delta\|h\|/\eta.
\]
Indeed $\|M'[h]v-\bar B_1h\|\le16\delta\|h\|/\eta$;
the true high and computed low normalized spectral blocks are separated
by more than $1/50$, and their overlap is at most $2^{14}\delta$.
The top normalized component also controls $L'[h]$, since
$u^TS_LM'[h]v=L'[h]/t_0$.

For the curvature comparison use the globally invertible matrices
\[
 A_K=I-K_L+uu^T,\qquad
 \widehat A_K=k_*I-N_K+w_Kw_K^T.
\]
Their norms of inverses are at most $2/\eta$ and $4/\eta$, respectively;
their difference is at most $2^{13}\delta/\eta$.
The inverse identity consequently bounds the difference of their reduced
resolvents by $2^{17}\delta/\eta^3$.
To evaluate the true response in \eqref{eq:genresponse}, subtract
$L'[h]v$ from $M'[h]v$. The resulting change of normalized forcing is
bounded by $2^{34}\delta\|h\|/\eta$, using the previous high-forcing
estimate. The computed forcing lies exactly in the computed low space,
so its reconstructed full reduced response is exactly $R_K$.
Using the bounded norms of both forcings and adding the direct
$\mathcal U-\mathcal D$ perturbation gives the safe contract
\begin{equation}\label{eq:genproxyerror}
 |L''[h,h]-h^T\bar Qh|
 \le2^{40}\delta\|h\|^2/\eta^3.
\end{equation}
With the unchanged $\delta=\eta^5/2^{100}$ this is at most
$\eta^2\|h\|^2/2^{60}$, well within the original curvature allowance.
The same precision also makes the trace and eigenvalue normalization
errors in the joint count, including $b_t\eta/(L-B_0)$, less than
$1/10000$ times $s$.

For the parameters used below, $\beta\|g\|\le25/8$ and
$\beta\|f''\|_\infty d^2\le9/4$. The normalized Gram norm is less than
$1+256\delta$ and $\|R_K\|<4$. Hence $\|\bar Q\|<128$.
The exact feasible negative basis is obtained from
$P(\bar Q-\eta I)P+(I-P)$ with diagonalization accuracy
$\eta/2^{24}$, exactly as in the joint-response construction.

\section{Finite comparison with normalized spectral cancellation}

\begin{lemma}\label{lem:generalizedperron}
Let $M(t)=B+C_0C_0^T+\sum_r a_re^{\ell_r(t)}z_rz_r^T$ be a PSD,
entrywise nonnegative Gram curve, with $B$ diagonal and
$0\preceq B\preceq(7/8)I$, $C_0,z_r\ge0$.
Assume the original logarithmic-curve bounds
$|\ell_r'(0)|\le K$, $-K^2/2\le\ell_r''\le0$, and
$|\ell_r'''|\le K^3/4$ almost everywhere, with $K\le1$.
At zero assume $\ell_0\le L\le9/8$, gap at least
$\zeta\in(0,1]$, and unit Perron vector $v>0$ with minimum $\nu$.
Set $k_v=1+\log(1/\nu)$ and form $K_L,S_L$ at this anchor.
Suppose
\[
 e=\|P_{(3/4,1]}(K_L)S_LM'(0)v\|
       \le K\zeta\nu/2^{16}.
\]
Then for $|t|\le\sqrt\zeta/(4096Kk_v)$ in the curve interval,
\[
 \Phi(t)\le\Phi(0)+\Phi'(0)t+\tfrac12\Phi''(0)t^2
        +2^{16}K^3|t|^3+2^{20}K^4t^4/\zeta,
\]
where $\Phi=(\lambda_1(M)-1)_+^3$ minus any quadratic energy.
\end{lemma}
\begin{proof}
Write $E=M'(0)$ and $a=v^TEv$. The positive Gram representation gives the
stronger entrywise estimate $|E|\le K(M-B)$ and therefore
\[
 |S_LEv|\le K D_L^{1/2}v=Kt_0u.
\]
Here $\sqrt{\ell_0-7/8}\le t_0\le\sqrt{9/8}<1.061$,
$\|S_L\|<3$, and $\min_i u_i\ge\nu/4$.
The $K_L$ low-response series has, relative to $u_i$, terms bounded by
both $(33/32)Kt_0$ and $Kt_0(3/4)^j/\min u_i$.
Summing the smaller estimates bounds this low response componentwise by
$18Kt_0k_vu_i$, and its norm by $4Kt_0$.
The high response has norm at most $2e/\zeta$ and componentwise relative
size at most $8e/(\zeta\nu)$.

The exact reduced response $z=(LI-M)^\dagger Ev$ is obtained by applying
$S_L(I-K_L)^\dagger$ to $S_L(Ev-av)$, and then subtracting the multiple
of $v$ that makes $z\perp v$.
Since $|a|\le t_0 e$, its correction has norm at most
$6t_0e/\zeta$ before the final multiplication by $S_L$.
Their componentwise bounds after multiplication by $S_L$, relative to
$v_i$, total at most $64e/(\zeta\nu)\le K/1024<K/100$; their norm
bounds are no larger. The displayed smallness assumption therefore
controls both kinds of error. The transformation $S_Lu=v/t_0$ preserves the
componentwise estimate relative to $v$, and the final gauge subtraction
adds at most the norm of the unprojected vector. It follows that
\[
 \|z\|\le16K,\qquad |z_i|\le64Kk_vv_i,
 \qquad |a|\le K/64.
\]

The remaining comparison is a direct expansion for the original $M(t)$,
not a differentiation of the normalized matrix. Put
$A=LI-M(0)$, $Z=\operatorname{diag}(z_i/v_i)$, and $D=I+tZ$.
Then
\[
 D((L+at)I-M(t))D=A+tG+R(t),\qquad
 G=ZA+AZ-E+aI,\qquad Gv=0.
\]
The positive Gram curve bounds give $\|Ev\|\le LK\le9K/8$,
$\|E\|\le2K$,
$\|M''(t)\|\le6K^2$, $\|M'''(t)\|\le14K^3$ almost everywhere.
Expansion yields $\|R(t)\|\le2^{14}K^2k_v^2t^2$:
the quadratic $ZAZ$ term is at most $8192K^2k_v^2$,
the two linear cross terms total less than $259K^2k_v^2$,
and the Taylor and cubic terms contribute less than $16K^2k_v^2$.
Off the diagonal $|G_{ij}|\le129Kk_vM_{ij}$, so the weighted Laplacian
identity gives $|y^TGy|\le129Kk_vy^TAy$.
For $u_0\perp D^{-1}v$, set $y=D^{-1}u_0\perp v$.
The step restriction implies $\|D-I\|\le1/64$, and the complementary
quadratic form is bounded below by
\[
 {1-129/4096-1/1024\over(1+1/64)^2}\zeta\|u_0\|^2
 >\tfrac12\zeta\|u_0\|^2.
\]
Thus $\lambda_2(M(t))\le L+at-\zeta/2$.

The positive trial vector $\psi=v+tz$ has norm at least one. Its residual
centered at $L+at$ is
$t^2(E-aI)z+(M(t)-M(0)-tE)\psi$, with normalized norm at most
$128K^2t^2$. Hence its Rayleigh quotient obeys
$|r-L-at|\le128K^2t^2$ and $r-\lambda_2(M(t))>\zeta/3$.
The Rayleigh residual inequality gives
$0\le\lambda_1(M(t))-r\le2^{16}K^4t^4/\zeta$.

For the Rayleigh Taylor expansion, $A_2=L''(0)/2$ has magnitude at most
$21K^2$, the numerator's cubic coefficient at most $608K^3$, and its
quartic coefficient at most $768K^4$. The division by
$1+t^2\|z\|^2$ changes these by at most $4K^3$ and $5376K^4$,
respectively. Using $K|t|\le1/4096$ and the $14K^3$ third-derivative bound
therefore gives
\[
 |r-L-at-\tfrac12L''(0)t^2|\le1024K^3|t|^3.
\]
Composition with $(u-1)_+^3$, whose second derivative is 6-Lipschitz,
costs less than $2^{16}K^3|t|^3$. All scalar arguments are below two,
so the replacement of $r$ by the top eigenvalue costs at most
$3\cdot2^{16}K^4t^4/\zeta<2^{20}K^4t^4/\zeta$.
\end{proof}

\paragraph{Remaining algorithm contracts.}
Keep the ridge and weight-precision schedules unchanged. Use coordinate
flatness $16\|h\|_\infty\le K$ and the dyadic test $9q/k$; with
$k>s/400$ this gives $K^2\le\min(1,14400q/s)$.
Put $c_*=2^{-56}$ in the original three dyadic step tests.
The new remainder divided by $\eta t^2$ is bounded by
$2^{16}c_*+2^{20}c_*^2<2^{-16}$.
The first derivative, normalized forcing, and curvature errors above are
absorbed by the original $\delta=\eta^5/2^{100}$, since $\eta\le2^{-40}$
and $\nu\ge\eta/2$.
The exact potential curvature remains at most
$-(19/10)\eta\|h\|^2$. The phase budget and $O(nJ^2q^2)$ update count
therefore remain valid. One extra symmetric diagonalization per update
changes only the absolute constant in ordinary cubic work per update.

For screened rows the frozen cold terms are placed in $C_0C_0^T$.
The same positive domination still holds, and the omitted cold-weight
perturbation remains within the existing error budget. Row-energy
control, the motion clock, the $1/4094$ high row-energy dimension charge,
and the two exact tangencies are unchanged.
\section{The final rational certificate: discrepancy below 37.54}
\label{sec:finalparameters}

Use the following constants in the retained-row walk:
\begin{equation}\label{eq:finalparameters}
\begin{gathered}
 R=\frac{7139}{5000},\qquad d=\frac{52}{173},\qquad
 \Delta=\frac{2704}{29929},\qquad c=\frac23,\qquad
 \beta=\frac{17}{25},\\
 A_0=\frac{25}{4},\qquad B_0=\frac{7463}{10000},\qquad
 a_0=\frac{5389}{2500},\qquad b=\frac{341613}{42500},\qquad
 H_0=\frac{94893}{5000}.
\end{gathered}
\end{equation}
Keep $\theta=1001/1000$, $L_*=257/256$, and
$\ell_0=999999/1000000$.  The even coordinate profile is
\[
 f(x)=\frac{519}{104}\log(3-2|x|)
                       +\frac{173}{52}(|x|-1).
\]
It is nonnegative and concave, is $C^{2,1}$ on the cube interval, and
vanishes at the endpoints.  Its height satisfies
$f(0)=2.15557478671874\ldots<a_0$.  The exact identities are
\[
 \frac{(1+d|f'(x)|)^2}{-f''(x)}=\frac{78}{173},\qquad
 \rho=\frac{\beta d}{c}=\frac{1326}{4325},\qquad
 \beta(b-a_0)=4.
\]
The last equality preserves the scale-decrease argument for the
$r_i^{-4}$ weights.  It has not been weakened by the optimization.

\subsection{Initialization and the joint dimension count}
The initialization exponent is
\[
 X=\beta(H_0/R-a_0-b)=2.10707849502731\ldots.
\]
The exact positive Taylor certificate
\[
 \sum_{j=0}^{47}\frac{X^j}{j!}>\frac{2A_0}{B_0R^2}
\]
has relative margin $0.000992324074770956\ldots$.
Moreover $\beta H_0/R>2$ and $-H_0/R+a_0+b<-3$.
These are sufficient for row entry and cleanup; earlier convenient
numerical bounds such as $-10$ or $-4$ are not required.

Use the reciprocal-response tangent at $t=1/4$:
\[
 \frac{1-\lambda}{1+\lambda}
 -\left(\frac{23}{25}-\frac{32}{25}\lambda\right)
 =\frac{32(\lambda-1/4)^2}{25(1+\lambda)}\ge0.
\]
Thus the joint-count coefficients are
\[
 a_t=\frac{25}{23},\qquad b_t=\frac{32}{23}.
\]
Set the computed normalized spectral cutoff to $18/25$ and retain the
true forcing cutoff $3/4$.  The tangent zero is
$23/32<18/25$.  In the perturbed joint-inertia lemma, the coefficient
of $\epsilon s/\lambda_*$ is $2a_t=50/23$.
The true-high/computed-low separation exceeds $0.029$, and the same
finite-error allowances remain valid; an intermediate overlap bound
$2^{14}\delta$ replaces $2^{13}\delta$.

Put
\[
 c_p=\frac{\theta^4}{A_0(1-\beta/100)}
     =0.161740798067016\ldots.
\]
The fully joint protected/large-row condition holds:
\[
 B_0R^2\left(c_p+\frac{b_t\Delta}{\ell_0-B_0}\right)
 =0.999892399877647\ldots<1.
\]
By the joint charge lemma, the combined loss for large rows, protected
pairs, high normalized spectral equations, and positive curvature is at
most
\[
 \frac1{R^2}+a_t\rho+b_t\Delta+c_p(L_*-B_0),
\]
apart from the reserved numerical allowance.  Stop at $s\le336$.
For every larger active size, including both exact tangencies and the
row-Gram high-space exclusion, the dimension margin is
\begin{align*}
 1-\frac1{R^2}-a_t\rho-b_t\Delta-c_p(L_*-B_0)
 -\frac1{10000}-\frac2{337}-\frac1{4094}
 &=0.00257600292021615\ldots\\
 &>\frac1{400}.
\end{align*}
All inequalities in this paragraph are checked with rational arithmetic.
The strict margins do not rely on rounded decimal displays.

\subsection{Terminal reserve and exact discrepancy}
The Pythagorean identity $52^2+165^2=173^2$ gives
\[
 A_\Delta=\frac{1+\sqrt{1-d^2}}d=\frac{13}{2},\qquad
 2RA_\Delta=13R=\frac{92807}{5000}.
\]
For $s\le336$, use conditional expectation with weight $19/6$:
\[
 \|A_V(\varepsilon-x)\|_2^2
 +\frac{19}{6}\|\varepsilon-x\|_2^2
 \le\frac{25}{6}s\le1400.
\]
This gives residual discrepancy at most $\sqrt{1400}$ and displacement
norm at most $\sqrt{8400/19}$.  The two required terminal inequalities are
exact rational comparisons:
\[
 \left(H_0+13R\right)^2>1400,\qquad
 (b+13)^2>\frac{8400}{19}.
\]
Indeed $b+13=21.0379529411765\ldots$.  For a still-large row, the current sum is zero and the residual bound
$\sqrt{1400}<C_*$ applies. For a medium row, the barriers and tracking give
\begin{align*}
 |(A\varepsilon)_i|
 &\le H_0-br_i+\sqrt{8400/19}\sqrt{q_i}
             +2A_\Delta(R-\sqrt{q_i})\\
 &\le C_*+
   \bigl(\sqrt{8400/19}-b-2A_\Delta\bigr)\sqrt{q_i}<C_*.
\end{align*}
Completed rows use the same deletion bound and their strict entry or
barrier inequality. For $n\le336$ the displayed quadratic rounding works
directly at $x=0$. Consequently
\[
 \|A\varepsilon\|_\infty<C_*:=H_0+13R
 =\frac{1877}{50}=37.54.
\]
Preprocessing discards rows of $\ell_1$ norm at most 37; hence
$m_0<n^2/37^2$.  All signing and rounding operations still use one sign
per original input column.

\subsection{Finite implementation contracts}
The new gradient satisfies $\|g_r\|\le3$, so use $g_r/3$ as the unit
row-gradient test in deterministic flattening.  Also
\[
 \beta\|g_r\|\le\frac{51}{25}<\frac{25}{8},\qquad
 \beta\|f''\|_\infty\Delta
 =\frac{5304}{4325}<\frac94.
\]
The proxy norm remains below 128 and the generalized finite comparison
continues to apply, since $B_0<7/8$.  Keep
$\eta=1/(2^{40}JS)$, $\delta=\eta^5/2^{100}$, negative-basis accuracy
$\eta/2^{24}$, and step cap $2^{-56}$.
Use the dyadic flatness test $9q/k$, giving
$16\|h\|_\infty\le K$ and
$K^2\le\min(1,14400q/s)$ when $k>s/400$.
The same $C^{2,1}$ logarithmic-curve bounds and
$2^{16}K^3|t|^3+2^{20}K^4t^4/\eta$ remainder therefore hold.
The ordinary arithmetic and update-count guarantees are unchanged.

The standard-library script \nolinkurl{certificates/verify_constant.py} checks these
scalar conditions exactly and prints the reproducible report
\nolinkurl{certificates/results/verify_constant.json}.  Its logarithm enclosure uses
$\log3=2\sum_{k\ge0}2^{-(2k+1)}/(2k+1)$ with an explicit rational tail;
the initialization exponential uses the positive Taylor lower bound above.

\section{Cubic row processing with ordinary arithmetic}
\label{sec:randomscreen}

The result of this section concerns row processing, not the entire
direction computation. It removes the retained-row Gram statistics from
the screened implementation. All matrices used to construct a feasible
negative basis are still charged at their actual cost.

\begin{proposition}[Randomized row processing]\label{prop:rowprocessing}
Consider a spectral walk with the following contracts: at a state with
$s$ active coordinates it computes an exactly feasible matrix $T$ with
$(3/4)I\preceq T^TT\preceq I$ and $k\ge c s$ columns, where $c>0$ is
absolute; all vectors in its span satisfy the required curvature
inequality; its row-gradient tests have bounded norm; and its finite-step
proof holds for both signs. Assume exact tangency $x_V^Th=0$, and the
phase and flatness conditions of the accompanying signing walk.
There is a randomized implementation preserving its discrepancy guarantee
which spends
\[
 O(mn)+\widetilde O(n^3)
\]
expected field operations on row records, row evaluation, screening,
and final verification. This cost includes all retained rows, however
tall the original matrix is. Construction of $T$, its constraint data,
and its spectral proxies is excluded and must be added separately.
Every returned signing is verified; the implementation restarts on a
failed run.
\end{proposition}

\subsection{A flat direction with a covariance bound}
Draw independent signs $\xi\in\{-1,1\}^k$ and form
$h=r_0T\xi/16$, where the rational $r_0>0$ obeys
$3/4\le r_0^2k\le1$. Reject only if a required flatness test fails.
For a test vector $a$ of norm at most one, the usual Rademacher
exponential bound gives
\[
 \Pr\{|r_0a^TT\xi|>3\sqrt{q/k}\}\le 2e^{-9q/2}.
\]
Choosing the existing $q$ large enough for the test list gives acceptance
probability at least $1/2$. Squared rational comparisons implement these
tests. The norm bounds hold for every sample, and conditioning on
acceptance gives
\begin{equation}\label{eq:randomcov}
 \frac1{32}\le\|h\|\le\frac1{16},\qquad
 \mathbb E[hh^T\mid\text{state, acceptance}]
 \preceq\frac{1}{128k}I\preceq\frac{C_h}{s}I.
\end{equation}
Here $C_h=1/(128c)$ suffices. This follows by bounding the unnormalized
accepted covariance by the unconditional covariance
$r_0^2TT^T/256$. No independence between the acceptance event and the
sample is assumed. The mean number of samples is at most two.

Choose one common dyadic step cap $\tau$ throughout a phase, using the
phase's worst allowed flatness constant. Since $S/2<s\le S$, this changes
the update count only by an absolute factor. Let $t_+,t_-$ be the actual
one-sided distances to the cube boundary, each capped at $\tau$. Move by
$d h$, where
\[
 d=t_+\quad\text{with probability }\frac{t_-}{t_++t_-},
 \qquad
 d=-t_-\quad\text{with probability }\frac{t_+}{t_++t_-}.
\]
Conditionally on the state and $h$, $\mathbb E d=0$ and
$\mathbb E d^2=t_+t_-$. To sample the rational probability, reveal fair
binary digits until their dyadic interval lies on one side of its
threshold; the expected number of comparisons is bounded by an absolute
constant. No random-real primitive is required.
Each short outcome freezes a coordinate. Tangency
gives the exact identity
$\|x+dh\|^2-\|x\|^2=d^2\|h\|^2$.
Consequently, over every run,
\begin{equation}\label{eq:opportunitybudget}
 \sum_{\text{accepted moves}}\tau^2\le1025n.
\end{equation}
Full outcomes cost at least $\tau^2/1024$ of coordinate energy; there are
at most $n$ short outcomes and $\tau\le1$. Thus the sum charges the
\emph{opportunity} to move by $\tau$, even when the realized move is short.
The original both-sign descent and phase proof still applies.

\subsection{Row energy in expectation}
Let $B$ denote the current matrix of retained coefficients, and let
$d_i=c_i+b_i^Tx_V$. These symbols are used for analysis; the algorithm
does not form $B^TB$. Cleanup preserves every $d_i$ and
$\operatorname{tr}(B^TB)\le s$. Conditional mean zero and
\eqref{eq:randomcov} imply
\[
 \mathbb E\left[\Delta\sum_i d_i^2\mid\text{state}\right]
 =\mathbb E[d^2 h^TB^TBh\mid\text{state}]
 \le\tau^2\operatorname{tr}\bigl(B^TB\,\mathbb E[hh^T]\bigr)
 \le C_h\tau^2.
\]
Starting at zero and using \eqref{eq:opportunitybudget} proves
\begin{equation}\label{eq:expectedrowenergy}
 \mathbb E\sum_i d_i^2\le1025C_h n
\end{equation}
at every stopped time in a bounded run. No row-energy tangency or
high-eigenvalue constraint on $B^TB$ is needed. This frees their direction
budget as well as their maintenance cost.

Use the same score $\chi_i=|d_i|+C_c r_i$ and hysteresis as in the
deterministic screening argument. The warm set $\mathcal W$ satisfies
pointwise
\[
 |\mathcal W|\le
 \frac{4\sum_i d_i^2+4C_c^2\theta^2s}{(H_0-2)^2}.
\]
Thus its expected size is $O(n+s\log^2 n)$, and direct evaluation of all
warm rows at all $\widetilde O(n)$ states costs
$\widetilde O(n^3)$ in expectation.

\subsection{A variance clock for sleeping rows}
Let $T_0=\widetilde O(n)$ be a deterministic cap for the number of accepted
moves from the finite-step proof; abort a run if a required dimension
certificate fails or this cap is reached. Also abort when a refreshed
barrier or computed spectral interval certifies violation of the admitted
range, before evaluating any exponential with an unbounded positive
argument. Use tolerances that cannot reject a state satisfying the proved
strict invariant. On every run, the matrices still have a known
polynomial norm bound, so the numerical work charged below remains valid
even after a missed cold certificate. Fix $\delta=1/8$ and choose
an integer $L_c$ such that
\[
 2e^{-L_c}\le\frac{\delta}{m_0(T_0+1)}.
\]
A sufficient rational construction replaces $e^{L_c}$ by $2^{L_c}$.
Here $m_0<n^2/C_0^2$ after discarding rows of $\ell_1$ norm at most a
fixed $C_0$ below the signing bound, so $L_c=O(\log(2+n))$.
Maintain the predictable clock
\[
 V=\frac1{256}\sum\tau^2.
\]
At a cold inspection with $\chi_i\le H_0-1$, save the current scale $r_i$
and issue a certificate lasting for clock increment
$q_i^{\rm clk}=1/(2L_c r_i^2)$.
Inspect before a proposed move would reach its expiration time. Impose
also $\tau^2/256\le1/(4L_cR^2)$; the phase step cap still has an
inverse-polylogarithmic lower bound.

While a certificate is active, the tracked row sum is a martingale with
increment bounded in absolute value by $r_i\tau/16$. Support changes
preserve the sum and only decrease its scale. The elementary
exponential-supermartingale bound therefore gives
\[
 \Pr\left\{\sup_{V-V_i\le q_i^{\rm clk}}|d_i-d_i(V_i)|\ge1
       \,\middle|\,\text{history at certification}\right\}
 \le2\exp\left(-\frac{1}{2r_i^2q_i^{\rm clk}}\right)
 =2e^{-L_c}.
\]
This follows from the conditional bound
$\mathbb E e^{\lambda\Delta d_i}\le
e^{\lambda^2r_i^2\tau^2/512}$ and the maximal inequality for the
resulting nonnegative supermartingale. It applies to adaptive stopping
times; it does not assume independent row sums.
There are at most $m_0(T_0+1)$ certificates. A union bound proves that all
cold score bounds are valid with probability at least $1-\delta$.
On that event the original small-weight and frozen-cold comparison
arguments apply at every move.

Successive due inspections of row $i$ are separated by at least
$1/(4L_c(r_i^{\rm ent})^2)$ in clock time, even when a warm period
intervenes. The factor two accounts for inspection before the next move;
the additional cap ensures that one clock increment uses at most half
a certificate's lifetime. Entry scales satisfy
$\sum_i(r_i^{\rm ent})^2\le\theta^2n$.
Since active dimension is nonincreasing and
\[
 \int s\,dV\le\frac{1025}{256}n^2,
\]
charging consecutive inspections to their disjoint clock intervals gives
total row-evaluation cost
\[
 O(m_0n)+O\left(L_c\sum_i(r_i^{\rm ent})^2\int s\,dV\right)
 =O(m_0n)+\widetilde O(n^3).
\]
Priority-queue operations are smaller. Classification at the single
entry of a row into the medium class costs another $O(m_0n)$.

\subsection{Support records, verification, and the remaining quartic term}
Store retained coefficients by original index, column incidence lists,
and a segment tree of retained squared coefficients for each row.
The tree supplies the least index above the current deletion threshold
in $O(\log n)$ work. Each freezing or deliberate deletion updates the
row's squared mass and offset in constant work, then its tree and scale
in logarithmic work. Every retained coefficient is removed at most once.
Removing a newly medium row's contribution from the large-row diagonal
also costs only its support size. Total support work is
$\widetilde O(m_0n)$; cold sums need not be re-evaluated at support changes.

Check the final signing against the original matrix in $O(mn)$ work.
Restart after a failed verification or an abort. On the event that all
clock certificates hold, the original correctness proof ensures success.
Thus each independent run succeeds with probability at least $7/8$;
the expected number of runs is at most $8/7$. All successful outputs are
explicitly verified. These observations prove
Proposition~\ref{prop:rowprocessing}.

The current construction of $T$ still forms dense coupled-response and
cancellation matrices and computes a negative subspace. At a state its
classical cost is $\widetilde O(w_s s^2+s^3)$, so its expected total cost
remains $\widetilde O(n^4)$. The preceding proposition removes a separate
row-processing obstacle; it supplies no amortization bound for these
changing dense subspaces. A cubic bound for the complete algorithm is
therefore not claimed.
\section{Complete modification and verification record}

For clarity, the deterministic algorithm proving Theorem~\ref{thm:main41}
uses the following changes to the fully specified algorithm in
\cite{baseline}. All omitted tie rules, support operations, and scalar
approximation procedures are unchanged.
\begin{enumerate}
\item Read the input once, discard rows of $\ell_1$ norm at most 37,
and use the profile and parameters in \eqref{eq:finalparameters}. If
$n\le336$, go directly to conditional-expectation rounding.
\item At each anchor refresh warm rows and expired deterministic cold
certificates. In the high branch compute both spectral approximations,
form the normalized high equations and proxy \eqref{eq:genproxy},
and construct the exact feasible negative basis at accuracy
$\eta/2^{24}$. Keep the two exact tangencies and the row-energy
high-space exclusion from the deterministic screened walk.
\item Flatten coordinates and the unit test vectors $(1-c)g_r$ for warm pairs. Use $h=h_0/16$,
the dyadic flatness test $9q/k$, and the largest dyadic $\tau$ satisfying
\[
 \tau\le2^{-56},\qquad K^2\tau\le2^{-56}\eta,
 \qquad K^2H^2\tau^2\le2^{-112}\eta,
 \quad H=\lceil\log_2(4/\eta)\rceil.
\]
Follow the original deterministic full-or-boundary rule, cleanup,
and phase resets. Stop at $s\le336$ and round with quadratic weight $\frac{19}{6}$.
\item For cold-weight screening set
$\eta_0=1/(2^{40}Jn)$, $\delta_0=\eta_0^5/2^{100}$ and
$\rho_c=\delta_0/[2^{22}(m_0+1)]$. Choose the least integer $\ell_c$
with $2^{\ell_c}\ge A_0/\rho_c$, and put
$C_c=a_0+b+\ell_c/\beta$. This replaces the old amplitude 36 by
$A_0=\frac{25}{4}$. All cold-weight errors retain the same bounds.
\end{enumerate}
The displayed dimension certificate includes every new equation.
The normalized finite comparison gives the same descent and phase
invariants, with the smaller fixed step cap. Cleanup and the remaining-mass
terminal estimate prove the signing bound. Each update still costs
$\widetilde O(w_s s^2+s^3)$ in ordinary arithmetic. The original row
screening gives $w_s=O(n+s\log^2 n)$, and the geometric phase sum proves
$O(mn)+\widetilde O(n^4)$ total work. This completes the theorem.

Section~\ref{sec:randomscreen} is an optional randomized implementation.
It replaces the deterministic flat-vector choice and boundary rule,
removes the two row-energy constraints, and uses a variance clock and
final verification. It retains the exact coordinate tangency. Its row
processing is cubic in expectation; its current complete implementation
still has quartic expected total work. The deterministic theorem does not
depend on that variant.

The package includes the standard-library exact parameter verifier
\texttt{verify\_constant.py} and its output. It certifies the logarithm
enclosure, initialization Taylor bound, joint deposit condition, full
dimension allowance, terminal inequalities, and finite-step scalar
constants. The new mathematical lemmas are proved above; numerical
experiments do not replace those proofs.

Two independent numerical audits test 300 examples of joint inertia,
480 identities for the normalized response, and 464 finite positive Gram
curves in total. The normalized response identity's largest relative error
in these tests was below $4\cdot10^{-15}$.
Exact rational checks additionally cover 670 joint budget examples,
404 supporting-line examples, the supporting-line identity,
large-row cancellation, covariance of accepted directions in a finite
example, unbiased boundary energy, and variance-clock accounting. The experiments
are basic checks, not end-to-end solver benchmarks or evidence of a cubic
runtime exponent. Selected parameter searches are archived with the certificates;
no numerical optimum is asserted to be global.

\begin{thebibliography}{9}
\bibitem{baseline} Project manuscript,
\emph{Constructive Koml\'os discrepancy below 67: coupled spectral signing
with input-linear row processing}, 27 September 2026,
\nolinkurl{archive/komlos_constructive_below67.tex} in this repository.
\bibitem{project} E. Akbas and S. Sra, project draft, 24 September 2026;
an extended version is \nolinkurl{archive/algo-komlos-development.tex} in
this repository; see also the algorithmic appendix of arXiv:2609.27172v1.
\bibitem{Li} X. Li, \emph{Fast Spectral Signing for Vector Balancing},
arXiv:2609.30044v1, 24 September 2026, Lemma 2.4 and Sections 4.3, 5, 9.
\url{https://arxiv.org/html/2609.30044v1}.
\end{thebibliography}
\end{document}
